\documentclass[UTF8,a4paper,12pt]{ctexart}

\usepackage{geometry}
\usepackage{setspace}
\usepackage{graphicx}
\usepackage{booktabs}
\usepackage{amsmath,amssymb}
\usepackage{caption}
\usepackage{fancyhdr}
\usepackage{titlesec}
\usepackage[hidelinks]{hyperref}

% Default report settings. Adjust these after auditing the official school template.
\geometry{
  left=2.5cm,
  right=2.5cm,
  top=2.5cm,
  bottom=2.5cm,
  headheight=14pt
}
\setstretch{1.5}
\setlength{\parindent}{2em}

% For strict school templates, set explicit Chinese and Latin fonts here.
% Example:
% \setCJKmainfont{SimSun}
% \setmainfont{Times New Roman}

\pagestyle{fancy}
\fancyhf{}
\fancyhead[C]{课程报告}
\fancyfoot[C]{\thepage}
\renewcommand{\headrulewidth}{0.4pt}
\renewcommand{\footrulewidth}{0pt}

\titleformat{\section}{\centering\heiti\zihao{-3}}{\thesection}{1em}{}
\titleformat{\subsection}{\heiti\zihao{4}}{\thesubsection}{1em}{}
\titleformat{\subsubsection}{\heiti\zihao{-4}}{\thesubsubsection}{1em}{}

\captionsetup{
  font=small,
  labelfont=bf,
  labelsep=space
}
\captionsetup[figure]{name=图}
\captionsetup[table]{name=表}

\newcommand{\schoolname}{学校名称}
\newcommand{\coursename}{课程名称}
\newcommand{\reporttitle}{报告题目}
\newcommand{\studentname}{姓名}
\newcommand{\studentid}{学号}
\newcommand{\teachername}{指导教师}
\newcommand{\submitdate}{\today}

\begin{document}

\begin{titlepage}
  \centering
  \vspace*{2cm}
  {\zihao{2}\heiti \schoolname \par}
  \vspace{2cm}
  {\zihao{1}\heiti \reporttitle \par}
  \vspace{2.5cm}
  {\zihao{4}
  \begin{tabular}{rl}
    课程： & \coursename \\
    姓名： & \studentname \\
    学号： & \studentid \\
    指导教师： & \teachername \\
    日期： & \submitdate \\
  \end{tabular}
  \par}
  \vfill
\end{titlepage}

\tableofcontents
\newpage

\section{引言}
在这里写研究背景、问题定义和报告结构。

\section{方法}
在这里写方法、实验设计或分析框架。

\subsection{数据与材料}
说明数据来源、实验环境或课程案例。

\section{结果与分析}
正文中引用图表时使用自动交叉引用，例如图~\ref{fig:example} 和表~\ref{tab:example}。

\begin{figure}[htbp]
  \centering
  \fbox{\rule{0pt}{4cm}\rule{0.8\linewidth}{0pt}}
  \caption{示例图片占位}
  \label{fig:example}
\end{figure}

\begin{table}[htbp]
  \centering
  \caption{示例表格}
  \label{tab:example}
  \begin{tabular}{lll}
    \toprule
    指标 & 结果 & 说明 \\
    \midrule
    A & 1.00 & 示例 \\
    B & 2.00 & 示例 \\
    \bottomrule
  \end{tabular}
\end{table}

\section{结论}
总结主要发现、局限性和后续工作。

\begin{thebibliography}{9}
\bibitem{sample} 作者. 文献题名[J]. 期刊名, 年份, 卷(期): 页码.
\end{thebibliography}

\end{document}
